\chapter{111}

\subsection{Louis}

Mein letztes Gespräch mit Joh hängt mir nach. Nach einigen Rückschlägen
glaube ich, einen Schritt weitergekommen zu sein. Ich mag seine
erklärenden Vergleiche. Gut, manchmal nicht. Ich will ehrlich sein,
beinahe immer dann nicht, wenn er ins Schwarze trifft. Wenn ich nicht
hören will, was er sagt.

«Könnte das wieder einer dieser Räuber sein, der dir reinpfuscht?» hatte
Joh gesagt und sich zurückgelehnt.

Dass er seinen Gedanken als Frage formuliert hat, hat bestimmt geholfen.
Überhaupt, er stellt mehrheitlich Fragen, fällt mir gerade auf, direkte
Aussagen macht er kaum.

Wenn ich mich ihm speziell nah fühle, ertappe ich mich beim Gedanken,
dass er auch ein Vater ist und ich mir wünsche, er wäre meiner. Gut,
vermag mich mein Selbstmitleid nur kurz ins Bein zu zwicken. Avas
Bildkommentare sind zu schön, um mich wegzudenken.

\sectionbreak

«Und für dich, Louis, hab ich das Fotobuch natürlich auch gemacht,» fügt
sie zum Schluss hinzu und blinzelt mich an, «vielleicht hilft es dir,
unsere \emph{Schafscheid}-Gespräche besser zu verstehen.»

Manuela grinst zu Joh hinüber. Er verhalten zurück. Was denken die sich
gerade? Ich kann nicht entscheiden, wer von uns beiden peinlicher
berührt ist, Ava oder ich. Ich muss vorsichtig sein. Will mir kein
zusätzliches Problem aufhalsen. Ich muss zugeben, es hat auch schon
geknistert zwischen uns. Sie ist eine tolle Frau. Doch damit hat es
sich. Ich will ihr keinesfalls irgendwelche Hoffnungen machen. Ich will
sie nicht verletzten. Dafür mag ich sie zu gern. Gut jedenfalls, dass
ich in den \emph{Boxenstopp} zu David gewechselt habe. Dass Ben zu Ava
in die \emph{Schäferloft} zog, hat sich geradezu aufgedrängt. Ich liess
mir nicht anmerken, wie dankbar ich für den Wechsel war. Ich glaube, wir
haben uns beide immer mal wieder schwer getan, den Schlafraum zu teilen.

Wie ich sie nun ansehe und ihr weiter zuhöre, werde ich unsicher.
Möglich, dass ich mir nur einbilde, ihr zu gefallen. Wie auch immer, es
ist gut, wenn ich mich künftig zurückhalte.

Auf Avas Handy ploppt eine Nachricht auf. Sie schaut kurz auf das
Display und springt hoch.

«Das ist Dad, bin gleich wieder da.»

Ich schaue ihr nach und höre ihre ersten Worte, bevor sie die Treppe
nach unten auf den Vorplatz verschwindet.

«Ja, natürlich, ich habe zwei machen lassen. Eins ist doch für dich,
Dad,» sagt sie und nach einer kurzen Pause und eine Spur leiser, «dann
sind die Bücher gleichzeitig angekommen, super. Und wie…»

Mehr verstehe ich nicht.

Die drei am Tisch übertreffen sich mit witzigen Kommentaren zu den
Fotos. Es scheint gleich weiterzugehen, wie gestern Abend. Sie sind ganz
aus dem Häuschen, so sehr freuen sie sich.

Dann kommt Avas Stimme wieder in Hörweite in die Küche zurück.

«Ja, klar, ich frage ihn. So gern geschehen, Dad, bis bald,» und damit
fixiert mich Ava.

«Louis, Dad fragt, ob du kurz mit ihm sprechen magst?»

«Ja,» ich zögere, damit habe ich nicht gerechnet, «doch, klar, gerne.»

Sie reicht mir ihr Handy und setzt sich wieder zu den andern an den
Tisch.

\sectionbreak

Erst draussen beginne ich mit Ben zu sprechen.

Nach Worten der Begrüssung kommt er auf den Punkt.

«Weisst du etwas Neues über den weiteren Verlauf der Ermittlungen,
Louis?»

«Ja, gerade gestern rief mich meine Mutter an. Sie ist ja in nahem
Kontakt zu Giorgias Eltern. Die haben ihre Informationen aus erster
Hand, also von der Polizei, denn über die Sache ist nach wie vor in
keiner Zeitung auch nur ein Wort zu lesen. Die Salas meinen, di Martino
habe es geschafft, sogar der Polizei einen Maulkorb zu verpassen.»

Ben unterbricht mich.

«Das spricht für seine Macht, hm, und zeigt dieses korrupte System.»

Ich nicke vor mich hin.

«Giorgias Anwalt soll eine Anklage wegen «Folgenschwerer Nötigung»
fordern. Maria hat sich bei der Polizei gemeldet, nachträglich, weil sie
es nicht mehr aushielt. Sie hat mitbekommen, dass Bastiano ihr ein
Getränk gegeben hat. Giorgia sei bald darauf komplett verändert gewesen.
Sie habe ihr Verhalten nicht verstanden. Die Polizei vermutet, dass er
da was reingemischt hat. Aber Bastiano streitet es ab.»

Ben fällt mir ins Wort. Die Verbindung ist nicht die beste.
Zwischendurch höre ich ihn abgehackt.

«Ich versteh ihren Anwalt. Ich wäre genauso vorgegangen. Doch, Chapeau,
der hat Courage, sich mit einem solchen Gegner anzulegen?»

«Ja, mir wurde ganz mulmig, als meine Mutter erzählte.»

«Hm, es fragt sich nun, wie taff Maria zu ihrer Aussage steht.»

«Wie meinst du das?»

«Wenn Aussage gegen Aussage steht, könnte es zu einer Gegenüberstellung
von Bastiano und Maria kommen. Er müsste dann auf ihre Aussage direkt
Stellung beziehen. Je nach Einschätzung der Polizei hätte schliesslich
das Gericht die Aufgabe, ein Urteil zu fällen.»

Bens Erklärungen lassen mich zusammensinken. Ich muss mich setzen. Die
Entwicklung der Geschichte macht mich fertig. Dieser hundserbärmliche
Scheissstock.

Auf Bens Frage nach meinem persönlichen Befinden kann ich ihm nicht viel
Neues erzählen.

«In Gesprächen mit Joh zeigt sich immer wieder, dass ich kaum
weiterkomme, bis ich Giorgia begegnet bin. Und davor habe ich immer noch
zu viel Angst.»

Kaum sind die Worte ausgesprochen, fühle ich mich ein bisschen besser.
Ich wage es, noch einen Schritt weiter zu gehen.

«Du weisst, das Schwierigste für mich ist einerseits nach wie vor, dass
sie sich nicht mehr an meine Anwesenheit erinnert, andererseits beruhigt
es mich ja auch, ach, Ben,» jetzt macht sich mein Hals doch wieder zu,
«ich, ich muss ihr mein verzweifeltes Handeln am Roccia dringend
erklären, nein, ich will es. Ich, ich komme sonst einfach nicht zu
meinem Frieden.»

«Ja, ich verstehe dich. Louis, du wirst genau spüren, wann der richtige
Zeitpunkt kommt, dich ihr zu stellen. Ich traue dir zu, dass du das
packst.»

Am liebsten würde ich Ben anspringen und umarmen. Seine Worte sind wie
eine dicke Schicht Balsam auf mein geschundenes Herz.
